


\section{Learning Credibility Scores On-The-Fly}\label{sec:tech-details}
 
Our system learns the credibility scores of the agents on the fly based on their performance in answering previous queries $\{q_1,\cdots, q_{t-1}\}$.

Initially, assuming there are no prior information about the reliability of the agents, all credibility scores are set to a default value (e.g., 0.5).
Then, at the end of each round $t$, the system computes a {\bf contribution score} $CSc^{(i)}$ for each of the team members $a_i\in A_t$.

Depending on the quality of the generated answer $o_t$ for the query $q_t$, the team is rewarded with a value $r_t\in[-1,1]$. The contribution scores and the reward value are then used for updating the credibility scores. The computation of the reward values is discussed in Appendix~\ref{app:reward}.

In the following, we first discuss the computation of the contribution scores, and then explain how the credibility scores are updated.

\subsection{Calculating the agent contributions}\label{sec:tech-details-Csc}

Given a query $q_t$, we model the process of generating the output $o_t$ as a game, where the team members collaboratively obtain the reward $r_t$.
Since the team members may have different impacts in the generation of the output, their share of the reward should be proportional to their contribution.

We propose the following approaches for computing the contribution scores (CSc):

\vspace{1mm}
\noindent{\em (i) Shapley Values for CSc computation:}
Our first approach for computing the contribution scores is based on {\em Shapley values} -- the popular concept in Game Theory for fairly distributing the reward among a team of players who have collaborated~\cite{shapley1951notes}.
Specifically, we consider Shapley values for {\em no-communication} topologies and when the aggregation strategy {\em is not} LLM-assisted.


Let $O=\{o_1,\cdots,o_N\}$ be the set of individual responses generated by a agents $A_t = \{a_1,\cdots,a_N\}$.
Let $\Sigma(S)$ be the final output generated by aggregating the responses of a subset of responses $S\subseteq O$.
Also, let $R(o_t)$ be the reward allocated based on the quality of $o_t$ as the answer of the query $q_t$.
The contribution score of the agent $a_i$ is then computed using the following equation:
\begin{align*}
 CSc^{(i)} =
\sum_{S\subseteq R\setminus{o_i}} \tfrac{|S|!(N-|S|-1)!}{N!}\Big(R&(\Sigma\big(S\cup{x_n})\big)\\
&-R\big(\Sigma(S)\big)\Big)
\end{align*}
%  A CSc value approaching $1.0$ indicates that an agent steered the team answer toward the optimum, whereas lower—or negative—scores signify that the agent pushed the solution further away from the correct outcome.

\vspace{-3mm}
\noindent{\em (ii) LLM-as-Judge for CSc computation:}
Despite their advantages such as theoretical guarantees, {\em it is \#P-hard} to compute Shapley values. As a result, computing the CSc values based on Shapley values require a {\em combinatorial} number of reward value computations for the aggregated outputs generated by each subset of $(A_t\backslash a_i)$.
This makes it practically infeasible to compute the contribution scores for the following cases.
(A) When the team members communicate, their output may be impacted by the composition of the team. As a result, for each subset $S\subseteq A_t$, one would need to form a new team and observe new outputs.
(B) When the aggregation of reward value computation is LLM-assisted, an LLM query would be needed for each subset $S\subseteq A_t$ to compute the reward.

Therefore, we instead use an LLM Judge to compute the contribution scores in such settings.
Specifically, given a query $q_t$,
we send the final answer $o_t$, the dialogue log, and the agent outputs to the LLM Judge, and ask the judge to quantify the contribution of each agent in the generation of the final output $o_t$.
The judge can analyze the message-passing log and observe which agents changed their response after the communication. The Agents never see these numbers to prevent strategic gaming.


\subsection{Updating the CrS values}  

Once the contribution scores are computed, the credibility score of each agent gets updated by distributing the reward $r_t$ among the agents proportional to their contribution. 
Specifically, using a learning rate $\eta$, the credibility scores are updated using Equation~\ref{eq:crs}.
\begin{equation}
\operatorname{CrS}^{(i)}_{t}=\operatorname{CrS}^{(i)}_{t-1}\bigl(1 + \eta . \operatorname{CSc}^{(i)}.r_t\bigr)
\label{eq:crs}
\end{equation}

Before concluding this section, we would like to remind that our scoring mechanism for computing CSc and CrS values is orthogonal to the team formation details including the agent roles and communication structure, making it easy to operate on top of the existing multi‑agent toolkits such as \textsc{AutoGen} and \textsc{CAMEL}.  Source code and prompts are provided in the supplementary material.